%=========================================================
\section{Soluciones}
%=========================================================

%=========================================================
\section{Ejercicios con números complejos}
%=========================================================

\subsection{Potencias de la unidad imaginaria}

Recordemos que
\[
i^2=-1.
\]

Por lo tanto, las primeras potencias de $i$ son
\[
i^1=i,\qquad
i^2=-1,\qquad
i^3=-i,\qquad
i^4=1.
\]

A partir de aquí, las potencias se repiten con periodo $4$:
\[
i,\,-1,\,-i,\,1,\,
i,\,-1,\,-i,\,1,\ldots
\]


\medskip

Para tener estas potencias como referencia, mostramos los valores desde
$n=0$ hasta $n=20$:

\[
\begin{array}{c|ccccccccccc}
n     &0&1&2&3&4&5&6&7&8&9&10\\ \hline
i^n   &1&i&-1&-i&1&i&-1&-i&1&i&-1
\end{array}
\]

\[
\begin{array}{c|cccccccccc}
n     &11&12&13&14&15&16&17&18&19&20\\ \hline
i^n   &-i&1&i&-1&-i&1&i&-1&-i&1
\end{array}
\]

La tabla permite observar nuevamente que el patrón se repite cada cuatro
potencias. En general, si
\[
n=4q+r,\qquad r\in\{0,1,2,3\},
\]
entonces
\[
\boxed{
i^n=
\begin{cases}
1, & r=0,\\
i, & r=1,\\
-1, & r=2,\\
-i, & r=3.
\end{cases}
}
\]

%---------------------------------------------------------
\subsubsection*{Ejemplo 1}

Calcular
\[
i^3.
\]

Utilizamos $i^2=-1$:
\[
i^3=i^2i=(-1)i.
\]

Por tanto,
\[
\boxed{i^3=-i}.
\]

%---------------------------------------------------------
\subsubsection*{Ejemplo 2}

Calcular
\[
i^{11}+i^{22}+i^{33}+i^{14}.
\]

Como las potencias de $i$ se repiten cada cuatro exponentes,
reducimos cada exponente módulo $4$.

Para $i^{11}$:
\[
11=4(2)+3,
\]
por lo que
\[
i^{11}=i^3=-i.
\]

Para $i^{22}$:
\[
22=4(5)+2,
\]
por lo que
\[
i^{22}=i^2=-1.
\]

Para $i^{33}$:
\[
33=4(8)+1,
\]
por lo que
\[
i^{33}=i.
\]

Finalmente, para $i^{14}$:
\[
14=4(3)+2,
\]
por lo que
\[
i^{14}=i^2=-1.
\]

Sustituyendo,
\[
i^{11}+i^{22}+i^{33}+i^{14}
=
-i-1+i-1.
\]

Agrupando las partes imaginarias y reales,
\[
(-i+i)+(-1-1)=0-2.
\]

Por tanto,
\[
\boxed{i^{11}+i^{22}+i^{33}+i^{14}=-2}.
\]

%=========================================================
\subsection{Evaluación de funciones trigonométricas}
%=========================================================

\subsubsection*{Ejemplo 3}

Calcular
\[
\cos\left(\frac{5\pi}{6}\right).
\]

Observamos que
\[
\frac{5\pi}{6}
=
\pi-\frac{\pi}{6}.
\]

Utilizamos la identidad
\[
\cos(\pi-\theta)=-\cos\theta.
\]

Entonces
\[
\cos\left(\frac{5\pi}{6}\right)
=
-\cos\left(\frac{\pi}{6}\right).
\]

Como
\[
\cos\left(\frac{\pi}{6}\right)=\frac{\sqrt{3}}{2},
\]
obtenemos
\[
\boxed{
\cos\left(\frac{5\pi}{6}\right)
=
-\frac{\sqrt{3}}{2}
}.
\]

%=========================================================
\subsection{Fórmula de Euler}
%=========================================================

Recordemos la fórmula de Euler:
\[
\boxed{e^{i\theta}=\cos\theta+i\sin\theta}.
\]

%---------------------------------------------------------
\subsubsection*{Ejemplo 4}

Evaluar
\[
e^{i\pi/6}.
\]

Aplicamos directamente la fórmula de Euler:
\[
e^{i\pi/6}
=
\cos\left(\frac{\pi}{6}\right)
+
i\sin\left(\frac{\pi}{6}\right).
\]

Sabemos que
\[
\cos\left(\frac{\pi}{6}\right)
=
\frac{\sqrt{3}}{2}
\]
y
\[
\sin\left(\frac{\pi}{6}\right)
=
\frac{1}{2}.
\]

Por lo tanto,
\[
e^{i\pi/6}
=
\frac{\sqrt{3}}{2}
+
i\frac{1}{2}.
\]

Así,
\[
\boxed{
e^{i\pi/6}
=
\frac{\sqrt{3}}{2}+\frac{i}{2}
}.
\]

%=========================================================
\subsection{Fórmula de De Moivre}
%=========================================================

En forma exponencial, la fórmula de De Moivre establece que
\[
\boxed{
\left(re^{i\theta}\right)^n
=
r^n e^{in\theta}
}.
\]

Equivalentemente,
\[
\left[r(\cos\theta+i\sin\theta)\right]^n
=
r^n\left(
\cos(n\theta)+i\sin(n\theta)
\right).
\]

%---------------------------------------------------------
\subsubsection*{Ejemplo 5}

Calcular
\[
\left(2e^{i\pi/5}\right)^7.
\]

Aplicamos De Moivre:
\[
\left(2e^{i\pi/5}\right)^7
=
2^7e^{i7\pi/5}.
\]

Como
\[
2^7=128,
\]
tenemos
\[
\left(2e^{i\pi/5}\right)^7
=
128e^{i7\pi/5}.
\]

Usando la fórmula de Euler,
\[
e^{i7\pi/5}
=
\cos\left(\frac{7\pi}{5}\right)
+
i\sin\left(\frac{7\pi}{5}\right).
\]

Por tanto,
\[
\boxed{
\left(2e^{i\pi/5}\right)^7
=
128\left[
\cos\left(\frac{7\pi}{5}\right)
+
i\sin\left(\frac{7\pi}{5}\right)
\right]
}.
\]

%=========================================================
\subsection{Cálculo exacto de $\cos(7\pi/5)$}
%=========================================================

Queremos determinar
\[
\cos\left(\frac{7\pi}{5}\right).
\]

Primero observamos que
\[
\frac{7\pi}{5}
=
\pi+\frac{2\pi}{5}.
\]

Utilizando
\[
\cos(\pi+\theta)=-\cos\theta,
\]
obtenemos
\[
\cos\left(\frac{7\pi}{5}\right)
=
-\cos\left(\frac{2\pi}{5}\right).
\]

Por lo tanto, necesitamos calcular exactamente
\[
\cos\left(\frac{2\pi}{5}\right).
\]

Sea
\[
\theta=\frac{2\pi}{5}
\]
y definamos
\[
z=e^{i\theta}=e^{i2\pi/5}.
\]

Entonces
\[
z^5=e^{i2\pi}=1.
\]

Así,
\[
z^5-1=0.
\]

Factorizando,
\[
(z-1)
\left(
z^4+z^3+z^2+z+1
\right)=0.
\]

Como $z\neq 1$, necesariamente
\[
z^4+z^3+z^2+z+1=0.
\]

Dividimos toda la ecuación entre $z^2$:
\[
z^2+z+1+\frac{1}{z}+\frac{1}{z^2}=0.
\]

Reagrupando,
\[
\left(z+\frac{1}{z}\right)
+
\left(z^2+\frac{1}{z^2}\right)
+1=0.
\]

Como
\[
z=e^{i\theta},
\qquad
\frac{1}{z}=e^{-i\theta},
\]
tenemos, por la fórmula de Euler,
\[
z+\frac{1}{z}
=
e^{i\theta}+e^{-i\theta}
=
2\cos\theta.
\]

Definamos
\[
x=2\cos\theta.
\]

Además,
\[
\left(z+\frac{1}{z}\right)^2
=
z^2+2+\frac{1}{z^2}.
\]

Por tanto,
\[
z^2+\frac{1}{z^2}
=
\left(z+\frac{1}{z}\right)^2-2
=
x^2-2.
\]

Sustituyendo en la ecuación anterior:
\[
x+(x^2-2)+1=0.
\]

Simplificando,
\[
x^2+x-1=0.
\]

Aplicamos la fórmula general:
\[
x
=
\frac{-1\pm\sqrt{1+4}}{2}
=
\frac{-1\pm\sqrt{5}}{2}.
\]

Como
\[
\theta=\frac{2\pi}{5}=72^\circ
\]
se encuentra en el primer cuadrante, sabemos que
\[
\cos\theta>0.
\]

Por tanto, elegimos
\[
x=\frac{-1+\sqrt{5}}{2}.
\]

Recordando que
\[
x=2\cos\theta,
\]
tenemos
\[
2\cos\theta
=
\frac{-1+\sqrt{5}}{2}.
\]

Dividiendo entre $2$:
\[
\cos\theta
=
\frac{\sqrt{5}-1}{4}.
\]

Así,
\[
\boxed{
\cos\left(\frac{2\pi}{5}\right)
=
\frac{\sqrt{5}-1}{4}
}.
\]

Finalmente,
\[
\cos\left(\frac{7\pi}{5}\right)
=
-\cos\left(\frac{2\pi}{5}\right),
\]
por lo que
\[
\boxed{
\cos\left(\frac{7\pi}{5}\right)
=
-\frac{\sqrt{5}-1}{4}
=
\frac{1-\sqrt{5}}{4}
}.
\]

%=========================================================
\subsection{Evaluación de $e^{i\pi/6}$}
%=========================================================

Evaluemos nuevamente
\[
e^{i\pi/6}.
\]

Por la fórmula de Euler,
\[
e^{i\theta}
=
\cos\theta+i\sin\theta.
\]

Tomando
\[
\theta=\frac{\pi}{6},
\]
resulta
\[
e^{i\pi/6}
=
\cos\left(\frac{\pi}{6}\right)
+
i\sin\left(\frac{\pi}{6}\right).
\]

Como
\[
\cos\left(\frac{\pi}{6}\right)
=
\frac{\sqrt{3}}{2},
\qquad
\sin\left(\frac{\pi}{6}\right)
=
\frac{1}{2},
\]
obtenemos
\[
e^{i\pi/6}
=
\frac{\sqrt{3}}{2}
+
\frac{1}{2}i.
\]

Por lo tanto,
\[
\boxed{
e^{i\pi/6}
=
\frac{\sqrt{3}}{2}
+
\frac{i}{2}
}.
\]

%=========================================================
\section{Ejercicios adicionales con números complejos}
%=========================================================

%---------------------------------------------------------
\subsection{Evaluación de una exponencial compleja}
%---------------------------------------------------------

\subsubsection*{Ejemplo 1}

Evaluar
\[
e^{1+i\pi}.
\]

Utilizamos la propiedad de la función exponencial
\[
e^{a+b}=e^a e^b.
\]

Por tanto,
\[
e^{1+i\pi}
=
e^1e^{i\pi}.
\]

Ahora aplicamos la fórmula de Euler:
\[
e^{i\theta}
=
\cos\theta+i\sin\theta.
\]

Para $\theta=\pi$,
\[
e^{i\pi}
=
\cos\pi+i\sin\pi.
\]

Como
\[
\cos\pi=-1,
\qquad
\sin\pi=0,
\]
tenemos
\[
e^{i\pi}
=
-1+i(0)
=
-1.
\]

Por consiguiente,
\[
e^{1+i\pi}
=
e(-1).
\]

Finalmente,
\[
\boxed{
e^{1+i\pi}=-e
}.
\]

En general, para $a,b\in\mathbb{R}$,
\[
\boxed{
e^{a+ib}
=
e^a(\cos b+i\sin b)
}.
\]

En efecto,
\[
e^{a+ib}
=
e^a e^{ib}
=
e^a(\cos b+i\sin b).
\]

%=========================================================
\subsection{Resolución de una ecuación cuadrática compleja}
%=========================================================

\subsubsection*{Ejemplo 2}

Resolver
\[
z^2-2z+5=0.
\]

Identificamos los coeficientes:
\[
a=1,
\qquad
b=-2,
\qquad
c=5.
\]

Aplicamos la fórmula general
\[
z
=
\frac{-b\pm\sqrt{b^2-4ac}}{2a}.
\]

Sustituyendo,
\[
z
=
\frac{-(-2)\pm
\sqrt{(-2)^2-4(1)(5)}}{2(1)}.
\]

Simplificando el discriminante,
\[
(-2)^2-4(1)(5)
=
4-20
=
-16.
\]

Por tanto,
\[
z
=
\frac{2\pm\sqrt{-16}}{2}.
\]

Recordemos que
\[
\sqrt{-16}
=
\sqrt{16}\sqrt{-1}
=
4i.
\]

Así,
\[
z
=
\frac{2\pm4i}{2}.
\]

Separando las dos posibilidades,
\[
z_1
=
\frac{2+4i}{2}
=
1+2i,
\]
y
\[
z_2
=
\frac{2-4i}{2}
=
1-2i.
\]

Por lo tanto, las soluciones son
\[
\boxed{
z_1=1+2i,
\qquad
z_2=1-2i
}.
\]

\medskip

\noindent
\textbf{Observación.}
Las raíces son complejos conjugados:
\[
z_2=\overline{z_1}.
\]

Esto es consecuencia de que el polinomio
\[
z^2-2z+5
\]
tiene coeficientes reales. En general, las raíces complejas no
reales de un polinomio con coeficientes reales aparecen en pares
conjugados.

%=========================================================
\subsection{Forma polar y exponencial de un número complejo}
%=========================================================

\subsubsection*{Ejemplo 3}

Expresar
\[
z=-1+\sqrt{3}\,i
\]
en forma polar y exponencial. Posteriormente calcular $z^{12}$.

\subsubsection*{Paso 1. Cálculo del módulo}

Si
\[
z=x+iy,
\]
su módulo está dado por
\[
|z|
=
\sqrt{x^2+y^2}.
\]

En nuestro caso,
\[
x=-1,
\qquad
y=\sqrt{3}.
\]

Por tanto,
\[
|z|
=
\sqrt{(-1)^2+(\sqrt{3})^2}.
\]

Entonces,
\[
|z|
=
\sqrt{1+3}
=
\sqrt{4}
=
2.
\]

Así,
\[
\boxed{|z|=2}.
\]

%---------------------------------------------------------
\subsubsection*{Paso 2. Determinación del argumento}

Tenemos
\[
x=-1<0
\]
y
\[
y=\sqrt{3}>0.
\]

Por consiguiente, el punto
\[
(-1,\sqrt{3})
\]
se encuentra en el \textbf{segundo cuadrante}.

Calculamos primero el ángulo de referencia $\alpha$.

Tenemos
\[
\tan\alpha
=
\left|
\frac{y}{x}
\right|.
\]

Por tanto,
\[
\tan\alpha
=
\left|
\frac{\sqrt{3}}{-1}
\right|
=
\sqrt{3}.
\]

Como
\[
\tan\left(\frac{\pi}{3}\right)
=
\sqrt{3},
\]
el ángulo de referencia es
\[
\alpha
=
\frac{\pi}{3}.
\]

Sin embargo, como $z$ se encuentra en el segundo cuadrante,
su argumento es
\[
\theta
=
\pi-\alpha.
\]

Por tanto,
\[
\theta
=
\pi-\frac{\pi}{3}
=
\frac{2\pi}{3}.
\]

Así, el argumento principal es
\[
\boxed{
\operatorname{Arg}(z)
=
\frac{2\pi}{3}
}.
\]

En general, todos los argumentos de $z$ están dados por
\[
\boxed{
\arg(z)
=
\frac{2\pi}{3}+2k\pi,
\qquad
k\in\mathbb{Z}
}.
\]

%---------------------------------------------------------
\subsubsection*{Paso 3. Forma polar}

La forma polar o trigonométrica de un número complejo es
\[
z
=
r(\cos\theta+i\sin\theta).
\]

Como
\[
r=2
\]
y
\[
\theta=\frac{2\pi}{3},
\]
obtenemos
\[
\boxed{
z
=
2\left(
\cos\frac{2\pi}{3}
+
i\sin\frac{2\pi}{3}
\right)
}.
\]

Podemos verificar el resultado utilizando
\[
\cos\left(\frac{2\pi}{3}\right)
=
-\frac{1}{2},
\qquad
\sin\left(\frac{2\pi}{3}\right)
=
\frac{\sqrt{3}}{2}.
\]

Entonces,
\[
2\left(
-\frac{1}{2}
+
i\frac{\sqrt{3}}{2}
\right)
=
-1+\sqrt{3}\,i,
\]
que es precisamente el número complejo original.

%---------------------------------------------------------
\subsubsection*{Paso 4. Forma exponencial}

Por la fórmula de Euler,
\[
e^{i\theta}
=
\cos\theta+i\sin\theta.
\]

Por tanto,
\[
r(\cos\theta+i\sin\theta)
=
re^{i\theta}.
\]

En nuestro caso,
\[
\boxed{
z
=
2e^{i2\pi/3}
}.
\]

%=========================================================
\subsection{Cálculo de una potencia mediante De Moivre}
%=========================================================

Ahora queremos calcular
\[
z^{12}.
\]

Como
\[
z
=
2e^{i2\pi/3},
\]
tenemos
\[
z^{12}
=
\left(
2e^{i2\pi/3}
\right)^{12}.
\]

Utilizamos la fórmula de De Moivre en forma exponencial:
\[
\boxed{
\left(re^{i\theta}\right)^n
=
r^n e^{in\theta}
}.
\]

Por tanto,
\[
z^{12}
=
2^{12}
e^{i12(2\pi/3)}.
\]

Calculamos por separado el módulo:
\[
2^{12}=4096.
\]

Para el argumento,
\[
12\left(\frac{2\pi}{3}\right)
=
\frac{24\pi}{3}
=
8\pi.
\]

Así,
\[
z^{12}
=
4096e^{i8\pi}.
\]

Aplicamos nuevamente la fórmula de Euler:
\[
e^{i8\pi}
=
\cos(8\pi)+i\sin(8\pi).
\]

Como $8\pi$ corresponde a cuatro vueltas completas,
\[
\cos(8\pi)=1
\]
y
\[
\sin(8\pi)=0.
\]

Por tanto,
\[
e^{i8\pi}
=
1.
\]

Finalmente,
\[
\boxed{
z^{12}=4096
}.
\]

\medskip

\noindent
\textbf{Observación.}
Este ejemplo muestra una de las principales ventajas de utilizar
la forma polar o exponencial de los números complejos. Calcular
directamente
\[
(-1+\sqrt{3}\,i)^{12}
\]
mediante productos algebraicos sucesivos sería considerablemente
más laborioso.

En cambio, al escribir
\[
-1+\sqrt{3}\,i
=
2e^{i2\pi/3},
\]
la potencia se obtiene simplemente elevando el módulo y
multiplicando el argumento por el exponente:
\[
r^n
\qquad\text{y}\qquad
n\theta.
\]

%=========================================================
\section{Soluciones adicionales del Primer Examen Parcial}
%=========================================================

En esta sección se incorporan únicamente los reactivos del examen que no
aparecen ya resueltos en las secciones anteriores. Los ejercicios que ya
cuentan con una solución desarrollada se conservan sin modificación.

%---------------------------------------------------------
\subsection{Potencias de $i$}
%---------------------------------------------------------

\subsubsection*{Problema 1}

Simplificar
\[
i^{11}+i^{22}+i^{33}+i^{44}.
\]

Como las potencias de $i$ se repiten con periodo $4$, reducimos cada
exponente módulo $4$.

Para $i^{11}$:
\[
11=4(2)+3,
\]
por lo que
\[
i^{11}=i^3=-i.
\]

Para $i^{22}$:
\[
22=4(5)+2,
\]
por lo que
\[
i^{22}=i^2=-1.
\]

Para $i^{33}$:
\[
33=4(8)+1,
\]
por lo que
\[
i^{33}=i^1=i.
\]

Finalmente, para $i^{44}$:
\[
44=4(11)+0,
\]
por lo que
\[
i^{44}=i^0=1.
\]

Sustituyendo estos valores en la expresión original,
\[
i^{11}+i^{22}+i^{33}+i^{44}
=
-i-1+i+1.
\]

Agrupamos las partes reales y las partes imaginarias:
\[
(-i+i)+(-1+1)=0+0.
\]

Por tanto,
\[
\boxed{
i^{11}+i^{22}+i^{33}+i^{44}=0
}.
\]

\medskip

\noindent
\textbf{Observación.}
La reducción módulo $4$ evita calcular sucesivamente todas las potencias de
$i$. Basta conocer el residuo del exponente al dividirlo entre $4$.

%---------------------------------------------------------
\subsection{Igualdad de números complejos}
%---------------------------------------------------------

\subsubsection*{Problema 2}

Determinar $a,b\in\mathbb{R}$ si
\[
(2a-1)+(3b+4)i=7-5i.
\]

Recordemos que dos números complejos
\[
z_1=x_1+y_1i
\qquad\text{y}\qquad
z_2=x_2+y_2i
\]
son iguales si y sólo si sus partes reales y sus partes imaginarias son
iguales, es decir,
\[
z_1=z_2
\quad\Longleftrightarrow\quad
x_1=x_2
\quad\text{y}\quad
y_1=y_2.
\]

En nuestro caso, la parte real del miembro izquierdo es
\[
2a-1,
\]
mientras que la parte real del miembro derecho es
\[
7.
\]

Por tanto,
\[
2a-1=7.
\]

Sumamos $1$ en ambos lados:
\[
2a=8,
\]
y dividimos entre $2$:
\[
\boxed{a=4}.
\]

Ahora igualamos las partes imaginarias:
\[
3b+4=-5.
\]

Restando $4$ en ambos lados,
\[
3b=-9.
\]

Dividiendo entre $3$,
\[
\boxed{b=-3}.
\]

Por consiguiente,
\[
\boxed{a=4,\qquad b=-3}.
\]

Podemos verificar el resultado sustituyendo estos valores:
\[
(2(4)-1)+(3(-3)+4)i
=
(8-1)+(-9+4)i
=
7-5i.
\]

Por tanto, los valores obtenidos satisfacen la igualdad original.

%---------------------------------------------------------
\subsection{Combinaciones lineales de números complejos}
%---------------------------------------------------------

\subsubsection*{Problema 3}

Sean
\[
z_1=3-2i,
\qquad
z_2=-1+5i,
\qquad
z_3=4+i.
\]
Calcular
\[
-5z_1+3z_2-4z_3.
\]

Sustituimos cada número complejo:
\[
-5(3-2i)+3(-1+5i)-4(4+i).
\]

Calculamos cada producto por separado.

Para el primer término,
\[
-5(3-2i)
=
-15+10i.
\]

Para el segundo término,
\[
3(-1+5i)
=
-3+15i.
\]

Para el tercer término,
\[
-4(4+i)
=
-16-4i.
\]

Por tanto,
\[
-5z_1+3z_2-4z_3
=
(-15+10i)+(-3+15i)+(-16-4i).
\]

Agrupamos las partes reales:
\[
-15-3-16=-34,
\]
y las partes imaginarias:
\[
10i+15i-4i=21i.
\]

Así,
\[
\boxed{
-5z_1+3z_2-4z_3=-34+21i
}.
\]

%---------------------------------------------------------
\subsection{Producto de números complejos}
%---------------------------------------------------------

\subsubsection*{Problema 4}

Calcular
\[
(2-i)(2+i)(1+3i).
\]

Observemos que los dos primeros factores son complejos conjugados:
\[
2-i
\qquad\text{y}\qquad
2+i.
\]

Podemos multiplicarlos directamente:
\[
(2-i)(2+i)
=
4+2i-2i-i^2.
\]

Los términos imaginarios se cancelan:
\[
2i-2i=0,
\]
y como
\[
i^2=-1,
\]
obtenemos
\[
-i^2=-(-1)=1.
\]

Por tanto,
\[
(2-i)(2+i)=4+1=5.
\]

La expresión original queda
\[
5(1+3i).
\]

Distribuyendo,
\[
5(1+3i)=5+15i.
\]

Finalmente,
\[
\boxed{
(2-i)(2+i)(1+3i)=5+15i
}.
\]

\medskip

\noindent
\textbf{Observación.}
En general, si $z=a+bi$, entonces
\[
z\overline z=(a+bi)(a-bi)=a^2+b^2=|z|^2.
\]
En este ejercicio,
\[
(2-i)(2+i)=|2+i|^2=5.
\]

%---------------------------------------------------------
\subsection{Conjugado de un producto}
%---------------------------------------------------------

\subsubsection*{Problema 5}

Verificar que
\[
\overline{zw}=\overline z\,\overline w
\]
para
\[
z=-3-i,
\qquad
w=3-i.
\]

Para verificar la identidad calcularemos ambos miembros por separado.

\subsubsection*{Cálculo de $\overline{zw}$}

Primero calculamos el producto $zw$:
\[
zw=(-3-i)(3-i).
\]

Desarrollando,
\[
(-3-i)(3-i)
=
(-3)(3)+(-3)(-i)+(-i)(3)+(-i)(-i).
\]

Calculamos cada término:
\[
(-3)(3)=-9,
\]
\[
(-3)(-i)=3i,
\]
\[
(-i)(3)=-3i,
\]
y
\[
(-i)(-i)=i^2=-1.
\]

Por tanto,
\[
zw=-9+3i-3i-1.
\]

Los términos imaginarios se cancelan:
\[
3i-3i=0,
\]
de modo que
\[
zw=-10.
\]

Como $-10$ es un número real, su conjugado es él mismo:
\[
\boxed{\overline{zw}=-10}.
\]

\subsubsection*{Cálculo de $\overline z\,\overline w$}

Los conjugados de $z$ y $w$ son
\[
\overline z=-3+i
\]
y
\[
\overline w=3+i.
\]

Por tanto,
\[
\overline z\,\overline w
=
(-3+i)(3+i).
\]

Desarrollando,
\[
(-3+i)(3+i)
=
-9-3i+3i+i^2.
\]

Como
\[
i^2=-1,
\]
obtenemos
\[
\overline z\,\overline w
=
-9-1
=
-10.
\]

Así,
\[
\boxed{\overline z\,\overline w=-10}.
\]

Como ambos miembros son iguales,
\[
\boxed{
\overline{zw}=\overline z\,\overline w
}.
\]

Por lo tanto, la propiedad queda verificada para los números complejos dados.

%---------------------------------------------------------
\subsection{Módulo de un número complejo}
%---------------------------------------------------------

\subsubsection*{Problema 6}

Calcular
\[
\left|1-\sqrt3\,i\right|.
\]

Recordemos que si
\[
z=a+bi,
\]
entonces su módulo se define como
\[
|z|=\sqrt{a^2+b^2}.
\]

En este caso,
\[
a=1
\qquad\text{y}\qquad
b=-\sqrt3.
\]

Por tanto,
\[
\left|1-\sqrt3\,i\right|
=
\sqrt{1^2+(-\sqrt3)^2}.
\]

Calculamos los cuadrados:
\[
1^2=1
\]
y
\[
(-\sqrt3)^2=3.
\]

Entonces,
\[
\left|1-\sqrt3\,i\right|
=
\sqrt{1+3}
=
\sqrt4.
\]

Finalmente,
\[
\boxed{
\left|1-\sqrt3\,i\right|=2
}.
\]

\medskip

\noindent
\textbf{Interpretación geométrica.}
El número complejo
\[
1-\sqrt3\,i
\]
corresponde al punto $(1,-\sqrt3)$ del plano complejo. Su módulo es la
distancia de este punto al origen:
\[
\sqrt{1^2+(-\sqrt3)^2}=2.
\]

%---------------------------------------------------------
\subsection{División de números complejos}
%---------------------------------------------------------

\subsubsection*{Problema 7}

Calcular
\[
\frac{3+2i}{1-i}.
\]

Para expresar el resultado en forma rectangular debemos eliminar el número
complejo del denominador. Para ello multiplicamos numerador y denominador
por el conjugado del denominador.

El conjugado de
\[
1-i
\]
es
\[
1+i.
\]

Por tanto,
\[
\frac{3+2i}{1-i}
\left(\frac{1+i}{1+i}\right)
=
\frac{(3+2i)(1+i)}{(1-i)(1+i)}.
\]

Calculamos primero el numerador:
\[
(3+2i)(1+i)
=
3+3i+2i+2i^2.
\]

Como
\[
i^2=-1,
\]
tenemos
\[
3+3i+2i-2
=
1+5i.
\]

Ahora calculamos el denominador:
\[
(1-i)(1+i)
=
1+i-i-i^2.
\]

Los términos imaginarios se cancelan y, usando $i^2=-1$,
\[
(1-i)(1+i)
=
1-(-1)
=
2.
\]

Por tanto,
\[
\frac{3+2i}{1-i}
=
\frac{1+5i}{2}.
\]

Separando las partes real e imaginaria,
\[
\boxed{
\frac{3+2i}{1-i}
=
\frac12+\frac52i
}.
\]

\medskip

\noindent
\textbf{Observación.}
La multiplicación por el conjugado funciona porque
\[
(a+bi)(a-bi)=a^2+b^2,
\]
que es siempre un número real.

%---------------------------------------------------------
\subsection{Conversión de forma polar a forma rectangular}
%---------------------------------------------------------

\subsubsection*{Problema 8}

Convertir
\[
5\left(
\cos\frac{5\pi}{6}
+i\sin\frac{5\pi}{6}
\right)
\]
a forma rectangular.

Recordemos que la forma polar o trigonométrica es
\[
z=r(\cos\theta+i\sin\theta),
\]
mientras que la forma rectangular es
\[
z=a+bi.
\]

En este ejercicio,
\[
r=5
\qquad\text{y}\qquad
\theta=\frac{5\pi}{6}.
\]

Observamos que
\[
\frac{5\pi}{6}
=
\pi-\frac{\pi}{6},
\]
por lo que el ángulo se encuentra en el segundo cuadrante.

En el segundo cuadrante el coseno es negativo y el seno es positivo.
Por tanto,
\[
\cos\left(\frac{5\pi}{6}\right)
=
-\cos\left(\frac{\pi}{6}\right)
=
-\frac{\sqrt3}{2},
\]
y
\[
\sin\left(\frac{5\pi}{6}\right)
=
\sin\left(\frac{\pi}{6}\right)
=
\frac12.
\]

Sustituimos estos valores:
\[
z
=
5\left(
-\frac{\sqrt3}{2}
+i\frac12
\right).
\]

Distribuyendo el factor $5$,
\[
z
=
-\frac{5\sqrt3}{2}
+
\frac52i.
\]

Por tanto, la forma rectangular es
\[
\boxed{
z=-\frac{5\sqrt3}{2}+\frac52i
}.
\]

%---------------------------------------------------------
\subsection{División en forma polar}
%---------------------------------------------------------

\subsubsection*{Problema 9}

Dividir
\[
6\operatorname{cis}\frac{5\pi}{6}
\]
entre
\[
3\operatorname{cis}\frac{\pi}{3}.
\]

Recordemos la notación
\[
\operatorname{cis}\theta
=
\cos\theta+i\sin\theta.
\]

Para dividir dos números complejos en forma polar utilizamos
\[
\frac{r_1\operatorname{cis}\theta_1}
     {r_2\operatorname{cis}\theta_2}
=
\frac{r_1}{r_2}
\operatorname{cis}(\theta_1-\theta_2).
\]

En nuestro caso,
\[
r_1=6,
\qquad
r_2=3,
\]
\[
\theta_1=\frac{5\pi}{6},
\qquad
\theta_2=\frac{\pi}{3}.
\]

Por tanto,
\[
\frac{6\operatorname{cis}(5\pi/6)}
     {3\operatorname{cis}(\pi/3)}
=
\frac63
\operatorname{cis}\left(
\frac{5\pi}{6}-\frac{\pi}{3}
\right).
\]

Calculamos el cociente de los módulos:
\[
\frac63=2.
\]

Ahora calculamos la diferencia de los argumentos. Escribimos
\[
\frac{\pi}{3}=\frac{2\pi}{6},
\]
de modo que
\[
\frac{5\pi}{6}-\frac{\pi}{3}
=
\frac{5\pi}{6}-\frac{2\pi}{6}
=
\frac{3\pi}{6}
=
\frac{\pi}{2}.
\]

Por tanto,
\[
\boxed{
\frac{6\operatorname{cis}(5\pi/6)}
     {3\operatorname{cis}(\pi/3)}
=
2\operatorname{cis}\frac{\pi}{2}
}.
\]

Si deseamos expresarlo en forma rectangular, utilizamos
\[
\operatorname{cis}\frac{\pi}{2}
=
\cos\frac{\pi}{2}
+i\sin\frac{\pi}{2}.
\]

Como
\[
\cos\frac{\pi}{2}=0
\qquad\text{y}\qquad
\sin\frac{\pi}{2}=1,
\]
tenemos
\[
2\operatorname{cis}\frac{\pi}{2}
=
2(0+i)
=
2i.
\]

Así, equivalentemente,
\[
\boxed{2i}.
\]

%---------------------------------------------------------
\subsection{Raíces sextas de la unidad}
%---------------------------------------------------------

\subsubsection*{Problema 12}

Encontrar las raíces sextas de la unidad.

Buscamos todos los números complejos $z$ que satisfacen
\[
z^6=1.
\]

La unidad puede escribirse en forma exponencial como
\[
1=e^{i2k\pi},
\qquad
k\in\mathbb Z,
\]
ya que
\[
\cos(2k\pi)=1
\qquad\text{y}\qquad
\sin(2k\pi)=0.
\]

Por tanto, la ecuación puede escribirse como
\[
z^6=e^{i2k\pi}.
\]

Si escribimos
\[
z=re^{i\theta},
\]
entonces
\[
z^6=r^6e^{i6\theta}.
\]

Comparando módulos,
\[
r^6=1.
\]
Como el módulo debe ser no negativo,
\[
r=1.
\]

Comparando argumentos,
\[
6\theta=2k\pi,
\]
de donde
\[
\theta=\frac{2k\pi}{6}
=
\frac{k\pi}{3}.
\]

Por tanto, las raíces están dadas por
\[
\boxed{
z_k=e^{ik\pi/3}
},
\qquad
k=0,1,2,3,4,5.
\]

Tomamos estos seis valores de $k$ porque a partir de $k=6$ las raíces
comienzan a repetirse.

Para $k=0$:
\[
z_0=e^{i0}=1.
\]

Para $k=1$:
\[
z_1=e^{i\pi/3}
=
\cos\frac{\pi}{3}
+i\sin\frac{\pi}{3}
=
\frac12+\frac{\sqrt3}{2}i.
\]

Para $k=2$:
\[
z_2=e^{i2\pi/3}
=
\cos\frac{2\pi}{3}
+i\sin\frac{2\pi}{3}
=
-\frac12+\frac{\sqrt3}{2}i.
\]

Para $k=3$:
\[
z_3=e^{i\pi}
=
-1.
\]

Para $k=4$:
\[
z_4=e^{i4\pi/3}
=
\cos\frac{4\pi}{3}
+i\sin\frac{4\pi}{3}
=
-\frac12-\frac{\sqrt3}{2}i.
\]

Para $k=5$:
\[
z_5=e^{i5\pi/3}
=
\cos\frac{5\pi}{3}
+i\sin\frac{5\pi}{3}
=
\frac12-\frac{\sqrt3}{2}i.
\]

Por consiguiente, las seis raíces sextas de la unidad son
\[
\boxed{
\left\{
1,
\frac12+\frac{\sqrt3}{2}i,
-\frac12+\frac{\sqrt3}{2}i,
-1,
-\frac12-\frac{\sqrt3}{2}i,
\frac12-\frac{\sqrt3}{2}i
\right\}
}.
\]

\medskip

\noindent
\textbf{Interpretación geométrica.}
Todas las raíces tienen módulo $1$, por lo que se encuentran sobre la
circunferencia unitaria. Además, sus argumentos consecutivos difieren en
\[
\frac{2\pi}{6}=\frac{\pi}{3}.
\]
Por ello, las seis raíces son los vértices de un hexágono regular inscrito
en la circunferencia unitaria.

